\chapter*{Заключение} \label{ch-conclusion}
\addcontentsline{toc}{chapter}{Заключение}	% в оглавление 

В рамках данной работы была решена задача разработки системы мониторинга производительности СУБД Greenplum. В результате проведённого исследования были проанализированы существующие подходы к мониторингу баз данных и выбраны современные инструменты, обеспечивающие сбор, хранение, визуализацию метрик и организацию системы оповещений.

Разработанная система реализована на основе экосистемы Prometheus и включает в себя полный набор компонентов, необходимых для комплексного мониторинга: сбор метрик с использованием экспортеров, их хранение и обработку, визуализацию в Grafana, а также механизм оповещений на базе Alertmanager. Все компоненты интегрированы в единую инфраструктуру с использованием контейнеризации Docker, что обеспечивает удобство развертывания и воспроизводимость системы.

В рамках практической части была проведена апробация разработанного решения с использованием тестовой базы данных TPC-H. В ходе тестирования выполнялись аналитические запросы, позволяющие создать нагрузку на систему и оценить её поведение. Результаты наблюдений показали, что разработанный дашборд в Grafana обеспечивает наглядное отображение изменений ключевых метрик, включая загрузку процессора, использование памяти и активность базы данных.

Дополнительно была проведена проверка системы оповещений. В результате искусственного создания критических ситуаций было подтверждено корректное срабатывание правил оповещений и доставка уведомлений пользователю. Это демонстрирует работоспособность и практическую применимость реализованного механизма контроля состояния системы.

Таким образом, разработанная система представляет собой не набор отдельных инструментов, а готовое интегрированное решение, обеспечивающее полный цикл мониторинга СУБД — от сбора и анализа метрик до оперативного оповещения о возникающих проблемах. Это подтверждает её практическую ценность и возможность использования в реальных условиях эксплуатации.

В качестве направлений дальнейшего развития можно выделить расширение набора метрик и правил оповещений, добавление новых каналов уведомлений, а также адаптацию системы для мониторинга других типов баз данных.